\documentclass{article}
\usepackage{amsmath, amssymb, amsthm, graphicx, xcolor}
\usepackage[margin=1in]{geometry}
\usepackage{hyperref}

% Theorem Environments
\newtheorem{theorem}{Theorem}
\newtheorem{definition}{Definition}
\newtheorem{lemma}{Lemma}
\newtheorem{example}{Example}

\title{\textbf{CSE400 - Fundamentals of Probability in Computing} \\ \Large Lecture 21 and L22: Inequalities and Tail Bounds}
\author{Dhaval Patel, PhD \\ \small Associate Professor \\ \small SEAS-Ahmedabad University, Ahmedabad, Gujarat, India}
\date{March 24 and 26, 2026}

\begin{document}

\maketitle

\tableofcontents
\newpage

\section{Tail and Motivation of Tail Bounds}

\subsection{Definition of a Tail}
\begin{definition}
The tail of a random variable $X$ is defined as the probability that $X$ exceeds a certain value $x$:
\[ P\{X > x\} \]
\end{definition}

\begin{figure}[h]
    \centering
    %     \textit{[Image Placeholder: A bell-shaped probability density curve. An arrow points from the center towards the right edge, labeled ``CALCULATE TAIL PROBABILITY''. A vertical line marks the value $x$ on the x-axis, and the area under the curve to the right of $x$ is shaded red and labeled ``Tail Area $P(X > x)$''.]}
    \caption{Probability Density and Tail Area}
\end{figure}

\subsection{Motivation: Why do we care about tails?}
Tails represent rare but significant events. We care about them because they often dictate system failures or worst-case scenarios.
\begin{itemize}
    \item Jobs delayed more than 24 hours in a scheduling system.
    \item Packets dropped due to a full buffer in a network.
\end{itemize}

\textbf{Key Idea:}
\begin{itemize}
    \item \textbf{Mean} $\rightarrow$ Average behavior (easy to compute and understand).
    \item \textbf{Tail} $\rightarrow$ Rare extreme events (hard to calculate).
\end{itemize}

\textbf{Why are tails hard to compute?}
They require summing many small probabilities, leading to complicated, long sums that are computationally intensive or analytically intractable.

\subsection{Tails Example 1: Job Distribution}
\begin{example}
\textbf{Q:} Suppose you are distributing $n$ jobs among $n$ servers at random. What's the probability that a particular server gets $\ge k$ jobs?

\textbf{Intuition:}
\begin{itemize}
    \item Jobs are assigned randomly.
    \item Most servers get few jobs.
    \item Occasionally, one server gets many jobs. This is a \textbf{tail event} (rare overload).
\end{itemize}

\textbf{Exact Computation:}
Let $X$ be the number of jobs assigned to a specific server. Then $X \sim \text{Binomial}(n, p)$, where $p = \frac{1}{n}$.
\[ P\{X \ge k\} = \sum_{i=k}^{n} \binom{n}{i} p^i (1-p)^{n-i} \]
This is a complicated, long sum!
\end{example}

\begin{figure}[h]
    \centering
    %     \textit{[Image Placeholder: Diagram showing jobs (colored circles) being randomly mapped via dashed arrows to bins representing servers $S_1$ to $S_n$. Server $n$ is highlighted in red, showing it has received an unusually high number of jobs ($\ge k$ jobs), illustrating a rare tail event.]}
    \caption{Random distribution of jobs leading to a tail event.}
\end{figure}

\subsection{Tails Example 2: Datacenter Arrivals}
\begin{example}
\textbf{Q:} Jobs arrive at a datacenter according to a Poisson process with rate $\lambda$ jobs/hour. What is the probability that $\ge k$ jobs arrive during the first hour?

\textbf{Intuition:}
\begin{itemize}
    \item Jobs arrive randomly over time.
    \item Most hours have \textit{moderate} arrivals.
    \item Sometimes, we see a \textbf{burst of arrivals}. This sudden spike in load is a \textbf{tail event}.
\end{itemize}

\textbf{Exact Computation:}
Let $X$ be the number of arrivals in one hour. Then $X \sim \text{Poisson}(\lambda)$.
\[ P\{X \ge k\} = \sum_{i=k}^{\infty} e^{-\lambda} \frac{\lambda^i}{i!} \]
Again, this requires computing a difficult infinite series.
\end{example}

\subsection{The Solution: Bounding Instead of Exact Computation}
Exact tail probabilities are hard to compute. \\
\textbf{Idea:} Instead of finding the exact value $P(X \ge k) = ?$, we find an \textbf{upper bound}.
\begin{itemize}
    \item It is easier to compute.
    \item It still gives a reliable \textbf{guarantee} for system performance.
\end{itemize}

\textbf{Key Idea:} Approximate safely instead of computing exactly:
\[ P(X \ge k) \le \text{something simple} \]

\subsubsection{Running Example for Bounds}
We will develop progressively better (tighter) tail bounds and test them on the same example:
\begin{itemize}
    \item \textbf{Experiment:} Flip a fair coin $n$ times.
    \item \textbf{Distribution:} $X = \text{number of heads}$. $X \sim \text{Binomial}(n, \frac{1}{2})$.
    \item \textbf{Mean:} $E[X] = \frac{n}{2}$.
    \item \textbf{Question:} What is $P(X \ge \frac{3}{4}n)$?
\end{itemize}

\begin{figure}[h]
    \centering
    \textit{[Image Placeholder: Bar chart of a Binomial(n, 1/2) distribution. The mean is marked at $n/2$. A threshold is marked at $3n/4$, and the bars to the right of this threshold are highlighted in red, indicating the tail region: "getting unusually many heads".]}
    \caption{Binomial distribution tail region.}
\end{figure}

\section{Markov's Inequality}

\subsection{Definition and Formal Statement}
Markov's inequality provides a basic upper bound for the tail probability of a non-negative random variable, using only its expected value.

\begin{theorem}[Markov's Inequality]
If $X$ is any non-negative random variable ($X \ge 0$) and $a > 0$, then:
\[ P(X \ge a) \le \frac{E[X]}{a} \]
\end{theorem}

\subsection{Proof of Markov's Inequality}
\begin{proof}
Let $X$ be a continuous non-negative random variable with probability density function $f(x)$. The expected value is:
\begin{align*}
E[X] &= \int_{0}^{\infty} x f(x) dx \\
&= \int_{0}^{a} x f(x) dx + \int_{a}^{\infty} x f(x) dx
\end{align*}
Since $X$ is non-negative, $x f(x) \ge 0$ for all $x$, so the first integral is non-negative:
\begin{align*}
E[X] &\ge \int_{a}^{\infty} x f(x) dx
\end{align*}
In the domain of the integral, $x \ge a$, which implies:
\begin{align*}
E[X] &\ge \int_{a}^{\infty} a f(x) dx \\
E[X] &\ge a \int_{a}^{\infty} f(x) dx \\
E[X] &\ge a P(X \ge a)
\end{align*}
Dividing both sides by $a$ yields:
\[ P(X \ge a) \le \frac{E[X]}{a} \]
\end{proof}

\subsection{Example}
Let's apply Markov's inequality to our running Binomial example.
\begin{itemize}
    \item $X \sim \text{Binomial}(n, \frac{1}{2})$
    \item $E[X] = \frac{n}{2}$
    \item We want to bound $P(X \ge \frac{3}{4}n)$.
\end{itemize}
Using Markov's Inequality ($a = \frac{3}{4}n$):
\[ P\left(X \ge \frac{3}{4}n\right) \le \frac{E[X]}{\frac{3}{4}n} = \frac{n/2}{3n/4} = \frac{2}{3} \]
\textbf{Conclusion:} Markov's inequality tells us the probability is at most $2/3$. This is a valid, but very loose (weak) bound.

\section{Chebyshev's Inequality}

\subsection{Definition and Formal Statement}
Chebyshev's inequality improves upon Markov's bound by utilizing both the mean \textit{and} the variance of the random variable, measuring how much the variable deviates from its mean.

\begin{theorem}[Chebyshev's Inequality]
For any random variable $X$ with expected value $\mu = E[X]$ and variance $\sigma^2 = \text{Var}(X)$, and for any real number $a > 0$:
\[ P(|X - \mu| \ge a) \le \frac{\sigma^2}{a^2} \]
\end{theorem}

\subsection{Proof of Chebyshev's Inequality}
\begin{proof}
We use Markov's Inequality. Define a new random variable $Y = (X - \mu)^2$.
Notice that $Y$ is strictly non-negative ($Y \ge 0$).
We apply Markov's inequality to $Y$ with the threshold $a^2$:
\[ P(Y \ge a^2) \le \frac{E[Y]}{a^2} \]
Substitute $Y = (X - \mu)^2$ back into the inequality:
\[ P((X - \mu)^2 \ge a^2) \le \frac{E[(X - \mu)^2]}{a^2} \]
Observe two facts:
1. The event $(X - \mu)^2 \ge a^2$ is exactly equivalent to the event $|X - \mu| \ge a$.
2. By definition, $E[(X - \mu)^2] = \text{Var}(X) = \sigma^2$.

Substituting these yields:
\[ P(|X - \mu| \ge a) \le \frac{\sigma^2}{a^2} \]
\end{proof}

\subsection{Example}
Applying Chebyshev to our running Binomial example:
\begin{itemize}
    \item $X \sim \text{Binomial}(n, \frac{1}{2})$
    \item $\mu = E[X] = \frac{n}{2}$
    \item $\sigma^2 = \text{Var}(X) = n p (1-p) = n(\frac{1}{2})(\frac{1}{2}) = \frac{n}{4}$
    \item We want to bound $P(X \ge \frac{3}{4}n)$.
\end{itemize}
Notice that if $X \ge \frac{3}{4}n$, then $X - \frac{n}{2} \ge \frac{1}{4}n$.
Thus, $P(X \ge \frac{3}{4}n) \le P(|X - \frac{n}{2}| \ge \frac{n}{4})$.
Using Chebyshev's Inequality with $a = \frac{n}{4}$:
\[ P\left(\left|X - \frac{n}{2}\right| \ge \frac{n}{4}\right) \le \frac{\sigma^2}{a^2} = \frac{n/4}{(n/4)^2} = \frac{n/4}{n^2/16} = \frac{4}{n} \]
\textbf{Conclusion:} Chebyshev gives us a bound of $4/n$. If $n=100$, the bound is $0.04$. This is \textit{much} tighter than Markov's $2/3$, because it decays as $n$ grows!

\section{Chernoff Bounds and Poisson Tail}

\subsection{Definition and MGF Derivation}
Chernoff bounds provide an exponentially tight bound on tail distributions by using the entire Moment Generating Function (MGF) instead of just the mean or variance.

\begin{definition}[Moment Generating Function]
The Moment Generating Function (MGF) of a random variable $X$ is defined as:
\[ M_X(t) = E[e^{tX}] \]
for any real number $t$.
\end{definition}

\begin{theorem}[Chernoff Bound]
For any random variable $X$ and any $t > 0$:
\[ P(X \ge a) \le e^{-ta} M_X(t) = \frac{E[e^{tX}]}{e^{ta}} \]
\end{theorem}

\subsection{Proof of Chernoff Bounds}
\begin{proof}
For any $t > 0$, the function $f(x) = e^{tx}$ is strictly increasing and strictly positive.
Therefore, the event $X \ge a$ is mathematically equivalent to the event $e^{tX} \ge e^{ta}$.
\[ P(X \ge a) = P(e^{tX} \ge e^{ta}) \]
Let $Y = e^{tX}$. Since $Y$ is a non-negative random variable, we can apply Markov's Inequality:
\[ P(Y \ge e^{ta}) \le \frac{E[Y]}{e^{ta}} \]
Substitute back $Y = e^{tX}$:
\[ P(X \ge a) \le \frac{E[e^{tX}]}{e^{ta}} = e^{-ta} M_X(t) \]
Because this holds for \textit{any} $t > 0$, the tightest bound is found by minimizing over all $t > 0$:
\[ P(X \ge a) \le \min_{t > 0} e^{-ta} M_X(t) \]
\end{proof}

\subsection{Upper and Lower Tail Bounds}
By analogous logic, for the lower tail (where $t > 0$):
\[ P(X \le a) = P(-X \ge -a) = P(e^{-tX} \ge e^{-ta}) \le \frac{E[e^{-tX}]}{e^{-ta}} \]
This allows symmetric analysis of both distribution tails.

\subsection{Application to Binomial}
Applying Chernoff to the sum of independent indicator variables (like our running Binomial example).
Let $X = \sum_{i=1}^n X_i$, where $X_i \sim \text{Bernoulli}(p)$. Let $\mu = E[X] = np$.
The standard multiplicative Chernoff bounds state:
\begin{itemize}
    \item \textbf{Upper Tail:} For any $\delta > 0$, $P(X \ge (1+\delta)\mu) \le \left( \frac{e^\delta}{(1+\delta)^{(1+\delta)}} \right)^\mu \le e^{-\frac{\delta^2 \mu}{2+\delta}}$
    \item \textbf{Lower Tail:} For any $0 < \delta < 1$, $P(X \le (1-\delta)\mu) \le \left( \frac{e^{-\delta}}{(1-\delta)^{(1-\delta)}} \right)^\mu \le e^{-\frac{\delta^2 \mu}{2}}$
\end{itemize}

Returning to our running example $P(X \ge \frac{3}{4}n)$, where $\mu = \frac{n}{2}$. We set $(1+\delta)\frac{n}{2} = \frac{3}{4}n \implies \delta = \frac{1}{2}$.
\[ P\left(X \ge \frac{3}{4}n\right) \le e^{-\frac{(1/2)^2 (n/2)}{2 + 1/2}} = e^{-\frac{n/8}{2.5}} = e^{-\frac{n}{20}} \]
\textbf{Conclusion:} Chernoff bounds show the probability decays \textit{exponentially} with $n$. This is profoundly tighter than both Markov (constant) and Chebyshev (polynomial $1/n$).

\subsection{Poisson Tail Bounds}
For a Poisson random variable $X \sim \text{Poisson}(\lambda)$, the MGF is:
\[ M_X(t) = E[e^{tX}] = \sum_{k=0}^{\infty} e^{tk} e^{-\lambda} \frac{\lambda^k}{k!} = e^{-\lambda} \sum_{k=0}^{\infty} \frac{(\lambda e^t)^k}{k!} = e^{-\lambda} e^{\lambda e^t} = e^{\lambda(e^t - 1)} \]
Applying the Chernoff Bound framework to $X$:
\[ P(X \ge x) \le \min_{t > 0} \frac{e^{\lambda(e^t - 1)}}{e^{tx}} \]
Taking the derivative with respect to $t$ and setting it to zero yields the optimal $t = \ln(x/\lambda)$. Plugging this back in provides the tightest exponential tail bound for Poisson processes (vital for analyzing traffic spikes in the datacenter example).

\section{Jensen's Inequality}

\subsection{Convex Functions: Definition and Properties}
\begin{definition}[Convex Function]
A real-valued function $f$ is convex on an interval if, for any two points $x_1, x_2$ in the interval and any $t \in [0,1]$:
\[ f(tx_1 + (1-t)x_2) \le t f(x_1) + (1-t) f(x_2) \]
Geometrically, the line segment between any two points on the graph of the function lies strictly above the graph.
Alternatively, if $f$ is twice differentiable, it is convex if $f''(x) \ge 0$ for all $x$.
\end{definition}

\begin{figure}[h]
    \centering
    \textit{[Image Placeholder: A convex upward-opening curve. A secant line connects two points $(x_1, f(x_1))$ and $(x_2, f(x_2))$. The curve dips strictly below the secant line, illustrating the definition of convexity.]}
    \caption{Visual representation of a convex function.}
\end{figure}

\subsection{Formal Statement}
\begin{theorem}[Jensen's Inequality]
If $X$ is a random variable and $f$ is a convex function, then:
\[ f(E[X]) \le E[f(X)] \]
If $f$ is strictly convex, equality holds if and only if $X$ is a constant (i.e., $X = E[X]$ with probability 1).
\end{theorem}

\subsection{Proof}
\begin{proof}
Let $\mu = E[X]$. Since $f$ is convex, there exists a tangent line at $x = \mu$ that sits completely below the function. Let the equation of this line be $L(x) = a(x - \mu) + f(\mu)$, where $a$ is the slope of the tangent.
By the property of convexity, for all $x$:
\[ f(x) \ge L(x) = a(x - \mu) + f(\mu) \]
Since this holds for all $x$, it holds for the random variable $X$:
\[ f(X) \ge a(X - \mu) + f(\mu) \]
Take the expected value of both sides:
\begin{align*}
E[f(X)] &\ge E[a(X - \mu) + f(\mu)] \\
E[f(X)] &\ge a E[X - \mu] + f(\mu)
\end{align*}
By linearity of expectation, $E[X - \mu] = E[X] - \mu = \mu - \mu = 0$.
Thus:
\[ E[f(X)] \ge a(0) + f(\mu) = f(E[X]) \]
\end{proof}

\subsection{Examples}
\begin{example}
Let $f(x) = x^2$. Since $f''(x) = 2 \ge 0$, $f$ is a convex function.
Applying Jensen's Inequality:
\[ (E[X])^2 \le E[X^2] \]
This inherently proves that Variance is non-negative: $\text{Var}(X) = E[X^2] - (E[X])^2 \ge 0$.
\end{example}

\begin{example}
Let $f(x) = e^x$. The exponential function is strictly convex ($f''(x) = e^x > 0$).
By Jensen's Inequality:
\[ e^{E[X]} \le E[e^X] \]
This relation often serves as a foundational step when establishing moment generating functions for bounding behavior.
\end{example}

\section{Case Study and System-Level Understanding}

\subsection{Interpretation of Bounds in Practice}
When building actual computing systems (e.g., allocating server space for datacenters, or buffer sizes in network routers), exact probabilities are secondary to \textbf{guaranteed bounds}. Engineers ask: "How many servers do I need so that the probability of overloading is less than $10^{-6}$?"

By using tail bounds:
\begin{itemize}
    \item We dimension systems based on worst-case bounds.
    \item We prevent cascading failures caused by rare, extreme tail events.
\end{itemize}

\subsection{Comparison of Bounds: Tightness vs. Information Needed}
We analyzed three major bounds on the running Binomial example:

1. \textbf{Markov's Inequality:}
   \begin{itemize}
       \item \textbf{Bound:} $\approx 2/3$ (Constant)
       \item \textbf{Tightness:} Very loose (weak).
       \item \textbf{Requirement:} Needs \textbf{only} the Mean ($E[X]$) and non-negativity.
   \end{itemize}

2. \textbf{Chebyshev's Inequality:}
   \begin{itemize}
       \item \textbf{Bound:} $4/n$ (Polynomial decay)
       \item \textbf{Tightness:} Moderate. Shrinks to 0 as $n \to \infty$.
       \item \textbf{Requirement:} Needs the Mean ($E[X]$) and the Variance ($\text{Var}(X)$).
   \end{itemize}

3. \textbf{Chernoff Bounds:}
   \begin{itemize}
       \item \textbf{Bound:} $\approx e^{-n/20}$ (Exponential decay)
       \item \textbf{Tightness:} Extremely tight. Shrinks incredibly fast.
       \item \textbf{Requirement:} Needs the full Moment Generating Function ($E[e^{tX}]$) and assumption of independent trials.
   \end{itemize}

\subsection{Real-World Meaning}
\textbf{The Trade-off:} The more information you know about the distribution (Mean $\rightarrow$ Variance $\rightarrow$ MGF), the tighter the bound you can construct. 
In datacenter provisioning:
\begin{itemize}
    \item If you only know average arrival rates, you must heavily over-provision (Markov).
    \item If you know the variance of the traffic, you can save money by provisioning less (Chebyshev).
    \item If you can model traffic perfectly as sum of independent variables, you can dimension the exact optimal buffer size, achieving high reliability at minimal cost (Chernoff).
\end{itemize}

\end{document}